\section{Give a seek its own working image}
Sequential decoding carries one accepted image into the next delta. A seek begins
with a frame number instead. The image currently displayed might belong to an
entirely different position, so it cannot serve as that delta's predecessor.
Implement \texttt{video\_seek} to choose an independent full frame and replay
only the records needed for the requested position.

Pass the immutable file base in \texttt{\$a0}, the unsigned target frame index in
\texttt{\$a1}, a four-byte-aligned destination image in \texttt{\$a2}, and a
96-byte, eight-byte-aligned native seek context in \texttt{\$a3}. The file and
its offset directory must come from a successful \texttt{video\_index} call.
Keep both unchanged while seeking. This interface uses that earlier whole-file
validation; it does not re-index a new file merely because its address is reused.
The context provides the accepted file count, frame count and period, plus separate
scratch image storage. The directory holds at least one native offset per frame.
The caller supplies valid disjoint file, directory, context, image, scratch and
private-stack windows. Retain the saved registers, frame pointer and stack pointer,
and keep the stack sixteen-byte aligned across nested calls.

\begin{table}[htbp]
  \centering
  \caption{Seek context. Each field occupies eight native bytes.}
  \label{tab:video-seek-context}
  \begin{tabular}{r p{1.4in} p{2.9in}}
    \toprule
    Offset & Field             & Meaning                                        \\
    \midrule
    0      & File bytes        & Input, admitted complete-file count.           \\
    8      & Directory         & Input pointer to aligned native offsets.       \\
    16     & Frame count       & Input, one through 256.                        \\
    24     & Full-frame period & Input, one through 64.                         \\
    32     & Image capacity    & Input, at least 16\,384 bytes.                 \\
    40     & Scratch address   & Input pointer, four-byte aligned.              \\
    48     & Scratch capacity  & Input, at least 16\,384 bytes.                 \\
    56     & Target            & Output, requested frame index.                 \\
    64     & Replay start      & Output, latest full frame at or before target. \\
    72     & Records replayed  & Output, target minus start plus one.           \\
    80--95 & Retained suffix   & Leave unchanged.                               \\
    \bottomrule
  \end{tabular}
\end{table}

Success returns zero in \texttt{\$v0} and 16\,384 in \texttt{\$v1}. Only
success copies the complete reconstructed image to the destination and publishes
context offsets 56 through 79. Failure returns one and zero without storing to
any destination-image or context byte, including previously successful results.
Scratch contents are unspecified on either return. Source and directory bytes
remain immutable. These rules make a failed seek harmless to the displayed frame,
even when a late replayed record fails after earlier records changed scratch.

Reject an out-of-range target, invalid file/frame/period counts, insufficient image
or scratch capacity, and misaligned native pointers before directory or file access.
The early context-alignment test must precede context loads. When consulting an
offset, require at least the file's 24-byte header before it and sixteen readable
bytes after it. A record window ends at the next directory offset, or at file end
for the final record. Reject reversed, short or oversized windows before passing
one to the record routine. These local bounds checks do not replace whole-file
validation or permit mutation of the indexed source.

The worked target is six. Ordinary moving data starts at full frame four; the
extra-full data starts at five. In either case, a successful seek produces the
same frame-six words. It publishes different replay starts and counts:

\begin{table}[htbp]
  \centering
  \caption{Two permitted encodings of the same requested pixels.}
  \begin{tabular}{l r r l}
    \toprule
    Input      & Target & Start & Applied record indices \\
    \midrule
    Moving     & 6      & 4     & 4, 5, 6                \\
    Extra full & 6      & 5     & 5, 6                   \\
    Moving     & 3      & 0     & 0, 1, 2, 3             \\
    Moving     & 0      & 0     & 0                      \\
    \bottomrule
  \end{tabular}
\end{table}

Attach this stage to the completed reader with
\texttt{python3 packages/034-lossless-video/construction/video-seek-construction/attach.py --source reader-work.asm
  --output seek-work.asm}. It retains all earlier instructions and names the previous
entry \texttt{reader\_example\_main}. The new caller loads \texttt{animation.mlv},
indexes it, then seeks. Generate the moving file with the earlier fixture command.
Use the seek preparation helper to choose a target or capacity without changing
instructions. A recipe naming extra-full requires that actual fixture as the file;
the grading backend supplies it automatically during the corresponding observation.
Output commands refuse existing files.

The illustration uses a 192-byte private frame, including a separate 48-byte record
context at offset 128. Free-writing work may use another layout. Unfinished regions
return development status four. After the image-copy region is filled, a later
unfinished publication region can return four with a changed image; completed
failure one never permits that behavior. Compare observations at their named
checkpoints before continuing, then check the completed callable contract.
